\begin{exo}{}
	On considère le plan muni d'un repère $\vOIJ$ orthonormé et les trois points $A$, $B$ et $C$ de coordonnées respectives $\coordl{-2}{1}$, $\coordl{0}{3}$ et $\coordl{3}{0}$.

	\begin{enumerate}
		\item \begin{enumerate}
			      \item Déterminer les coordonnées des vecteurs $\vec{AB}$ et $\vec{AC}$.
			      \item Déterminer les coordonnées du vecteur $\vec{AB} + \vec{AC}$.
		      \end{enumerate}
	\end{enumerate}
	On considère le point $D$ vérifiant la relation : \cmath{\vec{AB} + \vec{AC} = \vec{AD}}
	\begin{enumerate}[resume]
		\item \begin{enumerate}
			      \item En notant $(x_D\;; y_D)$ les coordonnées du point $D$, justifier qu'on a les deux égalités :
			            \cmath{\begin{cases} x_D + 2 = 7 \\ y_D - 1 = 1 \end{cases}}
			      \item En déduire les coordonnées du point $D$.
		      \end{enumerate}
	\end{enumerate}
\end{exo}
